\chapter*{Introducción}
\addchaptertocentry{Introducción}
\markboth{Introducción}{Introducción}

La tomografía computarizada (TC) es una técnica de diagnóstico por imágenes que permite reconstruir cortes transversales de un objeto a partir de mediciones externas de rayos X. Su desarrollo combina física, tecnología y matemática: las mediciones no entregan directamente la imagen interna, sino información indirecta obtenida a lo largo de trayectorias de rayos, lo que ubica la reconstrucción tomográfica dentro del marco de los problemas inversos \cite{kirsch2011}. En un modelo bidimensional idealizado, una sección del objeto se representa mediante una función que describe el coeficiente de atenuación, y los datos medidos se interpretan como integrales de esa función sobre rectas. Esta descripción conduce naturalmente a la transformada de Radon, introducida originalmente por Johann Radon en 1917 \cite{radon1917}.

Desde el punto de vista físico, el vínculo entre mediciones e integrales de línea proviene de la ley de Beer--Lambert. En un modelo ideal de rayos paralelos, al normalizar la intensidad medida y tomar logaritmos se obtiene una cantidad igual a la integral del coeficiente de atenuación sobre la trayectoria recorrida por el haz. Al repetir este proceso para distintas posiciones del detector y distintos ángulos de adquisición, los datos se organizan en un sinograma. Esta formulación constituye el punto de partida matemático de la reconstrucción tomográfica \cite{barrett_swindell1996,kuchment2014}.

El problema central consiste en reconstruir la función desconocida a partir de su transformada de Radon. En el modelo continuo, el teorema de corte de Fourier relaciona la transformada de Fourier de las proyecciones con valores de la transformada de Fourier bidimensional de la imagen. Esta relación permite derivar una fórmula de inversión mediante un cambio a coordenadas polares en frecuencia. El factor de densidad que aparece en ese cambio de variables da lugar al filtro rampa, cuyo símbolo es $|\sigma|$ bajo la convención de Fourier adoptada en esta tesis.

La retroproyección filtrada, o FBP por sus siglas en inglés, implementa esta idea en dos pasos: primero filtra cada proyección en la variable del detector mediante un multiplicador de Fourier y luego retroproyecta las proyecciones filtradas sobre el dominio espacial. La rampa clásica será la referencia teórica y computacional del método. Sin embargo, su crecimiento en altas frecuencias también explica la sensibilidad del procedimiento frente al ruido \cite[Cap.~3]{kak_slaney1988}\cite[Sec.~V.1]{natterer2001}; por ello, en la parte computacional se estudia una familia etiquetada A--F, explorada previamente en una contribución presentada en MAPI~4 \cite{velez_lopez2026mapi4}. Las definiciones matemáticas exactas empleadas en esta tesis se presentan en el Capítulo~3.

El diseño y análisis de filtros para FBP continúa siendo un tema activo, tanto desde el punto de vista de estimaciones de error como de la optimización de funciones filtro \cite{beckmann_iske2017,beckmann_nickel2025}. En esta tesis, los filtros A--F se usan con un propósito exploratorio y no reproducen los filtros optimizados de esos trabajos.

El objetivo de esta tesis es estudiar la retroproyección filtrada en reconstrucción tomográfica bidimensional desde una perspectiva teórica y computacional. En la parte teórica se fijan las convenciones de Fourier, se define la transformada de Radon, se deriva la retroproyección como adjunto formal y se obtiene la fórmula de FBP. En la parte computacional se analizan simulaciones reproducibles con un fantoma geométrico propio, sinogramas discretos y ruido controlado.

El alcance se limita a la geometría paralela bidimensional, al fantoma geométrico principal, al filtro rampa como referencia y a una familia A--F de filtros para comparación. No se pretende modelar un sistema clínico tridimensional ni establecer un filtro universalmente óptimo.

La estructura del trabajo es la siguiente. En el primer capítulo se presentan los preliminares estrictamente necesarios: convenciones de Fourier, nociones mínimas de operadores lineales y adjuntos, y multiplicadores de Fourier. En el segundo capítulo se desarrolla el modelo tomográfico continuo, la transformada de Radon, la retroproyección, el teorema de corte de Fourier y la fórmula de retroproyección filtrada. El tercer capítulo se dedica a la discretización, la implementación numérica y la comparación espectral y visual de filtros en reconstrucciones simuladas. El trabajo termina con conclusiones y líneas futuras.
